% ===== PRÉAMBULE COMMUN — À COPIER TEL QUEL EN TÊTE DE CHAQUE FICHIER .tex =====
% Ne modifier QUE :
%   - les deux lignes \fancyhead[L] et \fancyhead[R]
%   - le bloc titre \begin{center}...\end{center} juste après \begin{document}
\documentclass[11pt,a4paper]{article}
\usepackage[utf8]{inputenc}
\usepackage[T1]{fontenc}
\usepackage{lmodern}
\usepackage[french]{babel}
\usepackage{amsmath,amssymb,mathtools}
\usepackage{icomma}
\usepackage{xcolor}
\usepackage{tikz}
\usepackage{pgfplots}
\usepackage[most]{tcolorbox}
\usepackage{enumitem}
\usepackage{array,booktabs,colortbl}
\usepackage[a4paper,margin=2.2cm,headheight=26pt,headsep=10pt]{geometry}
\usepackage{fancyhdr}
\usepackage[hidelinks]{hyperref}
\usetikzlibrary{arrows.meta,calc,angles,quotes,positioning,decorations.markings,intersections,backgrounds}
\pgfplotsset{compat=1.18}

\definecolor{primary}{RGB}{226,74,88}
\definecolor{primarydark}{RGB}{150,28,42}
\definecolor{methcol}{RGB}{40,90,150}
\definecolor{excol}{RGB}{0,135,90}
\definecolor{defcol}{RGB}{0,114,178}
\definecolor{attncol}{RGB}{200,40,55}
\definecolor{thmcol}{RGB}{0,140,120}
\definecolor{courbe}{RGB}{32,110,200}
\definecolor{courbeder}{RGB}{210,70,40}
\definecolor{tang}{RGB}{0,150,90}
\definecolor{grille}{RGB}{200,205,215}
\definecolor{plus}{RGB}{200,60,60}
\definecolor{moins}{RGB}{30,90,200}

\tcbset{boxbase/.style={enhanced,breakable,boxrule=0.9pt,arc=2.5pt,
  left=3mm,right=3mm,top=2.3mm,bottom=2.3mm,fonttitle=\bfseries,coltitle=white}}
\newtcolorbox{cle}[1][]{boxbase,colback=defcol!6,colframe=defcol,
  title={\textsc{À retenir}\ifx\relax#1\relax\relax\else\textmd{ : \itshape #1}\fi}}
\newtcolorbox{methode}[1][]{boxbase,colback=methcol!7,colframe=methcol,sharp corners,
  title={\textsc{Méthode}\ifx\relax#1\relax\relax\else\textmd{ --- \itshape #1}\fi}}
\newtcolorbox{exemplebox}[1][]{boxbase,colback=excol!5,colframe=excol,
  title={\textsc{Exemple}\ifx\relax#1\relax\relax\else\textmd{ : \itshape #1}\fi}}
\newtcolorbox{piege}[1][]{boxbase,colback=attncol!6,colframe=attncol,
  title={\textsc{Attention}\ifx\relax#1\relax\relax\else\textmd{ : \itshape #1}\fi}}
\newtcolorbox{exo}[1][]{boxbase,colback=primary!4,colframe=primary,
  title={\textsc{Exercice}\ifx\relax#1\relax\relax\else\textmd{~#1}\fi}}
\newtcolorbox{corr}[1][]{boxbase,colback=thmcol!5,colframe=thmcol,
  title={\textsc{Corrigé}\ifx\relax#1\relax\relax\else\textmd{~#1}\fi}}

\newcommand{\R}{\mathbb{R}}
\newcommand{\N}{\mathbb{N}}
\newcommand{\ds}{\displaystyle}
\newcommand{\tg}{\operatorname{tg}}
\newcommand{\cotg}{\operatorname{cotg}}
\newcommand{\arctg}{\operatorname{arctg}}
\newcommand{\arccotg}{\operatorname{arccotg}}
\newcommand{\limh}{\lim_{h\to 0}}

\pagestyle{fancy}\fancyhf{}
\fancyhead[L]{\small\textcolor{primarydark}{\textbf{Thème 3} — Fonctions composées}}
\fancyhead[R]{\small\textcolor{primarydark}{\textsc{Tutoriel}}}
\fancyfoot[C]{\small\thepage}
\renewcommand{\headrulewidth}{0.8pt}

\begin{document}

\begin{center}
{\LARGE\bfseries\color{primarydark} Dériver une fonction composée}\\[2pt]
{\color{primary}\rule{0.5\textwidth}{1.2pt}}\\[4pt]
{\small\itshape La règle de la chaîne — les « pelures d'oignon »}
\end{center}
\vspace{2mm}

Une fonction \emph{composée} est une fonction \emph{emboîtée} dans une autre : $\sin(x^2)$, c'est
la fonction sinus \textbf{appliquée à} $x^2$. On l'imagine comme un \textbf{oignon} : plusieurs
\textbf{couches} superposées. Pour dériver, on épluche de l'extérieur vers l'intérieur.

\begin{cle}[le principe des pelures d'oignon]
On note $U$ la quantité \textbf{intérieure} (ce qui est « emballé »). On dérive la couche
\textbf{externe} en gardant $U$ intact, \textbf{puis} on multiplie par $U'$, la dérivée de l'intérieur.
Mini-formulaire :
\[
(U^n)'=nU^{\,n-1}U' \qquad (\sqrt{U})'=\frac{U'}{2\sqrt U} \qquad (e^{U})'=U'\,e^{U}
\qquad (a^{U})'=U'\,a^{U}\ln a
\]
\[
(\ln U)'=\frac{U'}{U} \qquad (\log_a U)'=\frac{U'}{U\ln a} \qquad
(\sin U)'=U'\cos U \qquad (\cos U)'=-\,U'\sin U
\]
\[
(\tg U)'=\frac{U'}{\cos^2 U} \qquad (\cotg U)'=-\frac{U'}{\sin^2 U} \qquad
(\arctg U)'=\frac{U'}{1+U^2}
\]
\[
(\arcsin U)'=\frac{U'}{\sqrt{1-U^2}} \qquad (\arccos U)'=-\frac{U'}{\sqrt{1-U^2}}
\]
Dans chaque formule, le facteur $U'$ est \textbf{la signature de la composée}.
\end{cle}

\begin{methode}[dériver une composée en 4 étapes]
\begin{enumerate}[leftmargin=6mm,itemsep=1pt]
\item \textbf{Repérer les couches} : quelle est la fonction \emph{externe} (celle du dessus) ?
      Que contient-elle ? On pose $U=$ l'intérieur.
\item \textbf{Dériver l'externe} avec la formule adaptée, \textbf{en gardant $U$ tel quel}.
\item \textbf{Multiplier par $U'$} (dérivée de l'intérieur). S'il reste des couches dans $U$,
      on recommence les étapes 1--3 sur $U$.
\item \textbf{Assembler et simplifier} le produit obtenu.
\end{enumerate}
\end{methode}

\begin{piege}[on oublie presque toujours le facteur $U'$ !]
C'est l'erreur numéro un. On dérive l'enveloppe mais on oublie de multiplier par la dérivée
de l'intérieur.
\[
(\sin x^2)' = 2x\cos x^2 \quad\text{et non}\quad \cos x^2 ;\qquad
(f^2)' = 2f'f \quad\text{et non}\quad 2f.
\]
Réflexe : après avoir dérivé la couche externe, demande-toi toujours
« \emph{et l'intérieur, quelle est sa dérivée ?} » — puis multiplie.
\end{piege}

\begin{exemplebox}[$y=\sin(x^2)$]
Couche externe : le sinus. Intérieur : $U=x^2$, donc $U'=2x$.\\
Formule $(\sin U)'=U'\cos U$ :
\[
y' = \underbrace{2x}_{U'}\cdot \underbrace{\cos(x^2)}_{\cos U} = 2x\cos(x^2).
\]
\end{exemplebox}

\begin{exemplebox}[$y=\sqrt{x^2+1}$]
Couche externe : la racine. Intérieur : $U=x^2+1$, donc $U'=2x$.\\
Formule $(\sqrt U)'=\dfrac{U'}{2\sqrt U}$ :
\[
y' = \frac{2x}{2\sqrt{x^2+1}} = \frac{x}{\sqrt{x^2+1}}.
\]
\end{exemplebox}

\begin{center}
\begin{tikzpicture}[scale=1]
  \foreach \i/\c/\lab in {0/defcol!12/{}, 1/methcol!14/{}, 2/excol!16/{}} {}
  % pile de couches (oignon)
  \draw[fill=defcol!10,draw=defcol,thick] (0,0) ellipse (3.1 and 0.85);
  \node[defcol] at (0,-0.03) {\small couche externe : $\sin(\ \cdot\ )$};
  \draw[fill=methcol!12,draw=methcol,thick] (0,0) ++(0,0) ellipse (2.0 and 0.55);
  \node[methcol] at (0,0) {\small intérieur : $U=x^2$};
  \draw[-{Stealth[length=3mm]},thick,primarydark] (3.4,0.4) to[out=30,in=90] (3.9,-0.1);
  \node[primarydark,right,align=left] at (3.7,-0.55)
     {\small dériver l'externe,\\[-1pt]\small puis $\times\,U'=2x$};
\end{tikzpicture}
\end{center}

\end{document}
